\appendixitem{ZARISKI TOPOLOGY}

Let V be a finite dimensional vector space. We denote by \(A_{V}\) the algebra of polynomial functions on V with values in k (Algebra, Chap. IV, §5, no. 10, Def. 4). This is a graded algebra; its component of degree 1 is the dual \(V^{*}\) of V, and the injection of \(V^{*}\) into \(A_{V}\) extends to an isomorphism from the symmetric algebra \(\mathbf{S}(V^{*})\) to \(A_{V}\) (Algebra, Chap. IV, §5, no. 11, Remark 2).

If \((e_{1},\ldots,e_{n})\) is a basis of V, and \((\mathrm{X}_{1},\ldots,\mathrm{X}_{n})\) a sequence of indeterminates, the map from \(k[\mathrm{X}_{1},\ldots,\mathrm{X}_{n}]\) to \(A_{V}\) that takes any element f of \(k[\mathrm{X}_{1},\ldots,\mathrm{X}_{n}]\) to the function

\[
\sum_ {i = 1} ^ {n} \lambda_ {i} e _ {i} \mapsto f (\lambda_ {1}, \ldots , \lambda_ {n})
\]

is an isomorphism of algebras (Algebra, Chap. IV, §5, no. 10, Cor. of Prop. 19).

PROPOSITION 1. Let H be the set of algebra homomorphisms from \(A_{V}\) to k. For any \(x \in V\) , let \(h_{x}\) be the homomorphism \(f \mapsto f(x)\) from \(A_{V}\) to k. Then, the map \(x \mapsto h_{x}\) is a bijection from V to H.

Indeed, let \(\mathrm{H}'\) be the set of algebra homomorphisms from \(k[\mathrm{X}_1,\ldots ,\mathrm{X}_n]\) to \(k\) . The map \(\chi \mapsto (\chi (\mathrm{X}_1),\dots ,\chi (\mathrm{X}_n))\) is clearly a bijection from \(\mathrm{H}'\) to \(k^n\) .

COROLLARY. For any \(x \in V\) , let \(\mathfrak{m}_{x} = \operatorname{Ker}(h_{x})\) . Then the map \(x \mapsto \mathfrak{m}_{x}\) is a bijection from V to the set of ideals m of \(A_{V}\) such that \(A_{V}/m = k\) .

A subset \(\mathrm{F}\) of \(\mathrm{V}\) is said to be closed if there exists a family \((f_i)_{i\in \mathrm{I}}\) of elements of \(\mathrm{A_V}\) such that

\[
x \in \mathrm{F} \Longleftrightarrow x \in \mathrm{V} \text {   and   } f _ {i} (x) = 0 \text {   for   all   } i \in \mathrm{I}.
\]

It is clear that \(\varnothing\) and V are closed, and that any intersection of closed sets is closed. If F is defined by the vanishing of the \(f_{i}\) and \(\mathrm{F}'\) by that of the \(f_j'\) , \(\mathrm{F} \cup \mathrm{F}'\) is defined by the vanishing of the \(f_{i}f_{j}'\) , and hence is closed. Thus, there exists a topology on V such that the closed sets for this topology are exactly the closed sets in the above sense. This topology is called the Zariski topology on V. For any \(f \in \mathrm{A_V}\) , we denote by \(\mathrm{V}_f\) the set of \(x \in \mathrm{V}\) such that \(f(x) \neq 0\) ; this is an open subset of V. It is clear that the \(\mathrm{V}_f\) form a base of the Zariski topology. (If \(k\) is a topological field, the canonical topology of V is finer than the Zariski topology.)

The map \(x \mapsto m_{x}\) of the Cor. of Prop. 1 can be considered as a map \(\varepsilon\) from V to the prime spectrum \(\operatorname{Spec}(A_{\mathrm{V}})\) of \(A_{V}\) (Commutative Algebra, Chap. II, §4, no. 3, Def. 4). It is immediate that the Zariski topology is the inverse image under \(\varepsilon\) of the topology of \(\operatorname{Spec}(A_{\mathrm{V}})\) .

PROPOSITION 2. The vector space V, equipped with the Zariski topology, is an irreducible noetherian space. In particular, every non-empty open subset of V is dense.

Since \(A_{V}\) is noetherian, \(\operatorname{Spec}(A_{V})\) is noetherian (Commutative Algebra, Chap. II, §4, no. 3, Cor. 7 of Prop. 11), and every subspace of a noetherian space is noetherian (loc. cit., no. 2, Prop. 8). With the notation of the Cor. of Prop. 1, the intersection of the \(m_{x}\) is \(\{0\}\) , and \(\{0\}\) is a prime ideal of \(A_{V}\) ; thus V is irreducible (loc. cit., no. 3, Prop. 14).

